% Job posting: https://ca.linkedin.com/jobs/view/sensor-fusion-software-engineer-at-hexagon-autonomous-solutions-4467284436
\documentclass[11pt]{article}
\usepackage[left=1in,top=1in,right=1in,bottom=1in]{geometry}
\usepackage[hidelinks]{hyperref}
\usepackage{setspace}
\usepackage[T1]{fontenc}
\usepackage{newtxtext,newtxmath}
\ifdefined\pdfgentounicode
  \input{glyphtounicode}
  \pdfgentounicode=1
\fi
\setstretch{1.1}
\pagenumbering{gobble}
\begin{document}
\pagestyle{empty}

\begin{flushright}
Nishal Stanislaus Silva, PhD \\
Toronto, ON \\
nishal.silva@hotmail.com \\
+1 613 200 2030
\end{flushright}

Dear Hiring Manager,

At McGill University's IDMIL/CIRMMT lab, I built a self-powered smart electric guitar that fused two independent real-time sensor streams -- an inertial (IMU) channel and an audio channel -- through a hierarchical finite-state machine driving a single-LSTM-256 gesture model, reaching F1 0.78 on a 500-event corpus. The engineering problem underneath that result is the one your Sensor Fusion Software Engineer role solves at a different scale: combining multiple imperfect, asynchronous sensor streams into a single trustworthy, real-time state estimate, running on embedded hardware with hard power and compute limits. Autonomous heavy-equipment vehicles fuse GPS, IMU, lidar, and radar for positioning and perception; I fused IMU and audio for gesture and pattern recognition -- the architecture, the arbitration logic, and the real-time constraints are the same class of problem, just applied to a different sensor suite.

That project sits inside a broader body of applied research built on the same discipline. As a postdoctoral researcher at the University of Trento, on the EU Horizon EIC Pathfinder project MUSMET, I owned a full ML pipeline for streaming audio and IMU/gesture data end to end -- acquisition, feature engineering, training, and edge deployment on Raspberry Pi 4 -- and designed frozen-protocol benchmark and ablation suites to quantify accuracy, latency, and CPU trade-offs before anything shipped. The resulting real-time pattern detector ran at 14 ms inference with F1 0.76, 74x faster than the 1038 ms baseline it replaced. I also built a real-time multi-sensor synchronization pipeline fusing audio, BCI/EEG, and mixed-reality head/hand tracking across four simultaneous performers, which demanded the same synchronization and power-profiling discipline your embedded sensor stack requires. Earlier, 5.5 years as a Research Engineer at MAS Holdings gave me the C++ systems-engineering foundation to take that kind of work from prototype to production: end-to-end computer-vision and IoT systems running on embedded hardware at manufacturing scale, with CI/CD and code review practices I introduced along the way.

I have not worked directly with GPS, lidar, or radar, and I have no mining or autonomous-vehicle industry experience specifically -- I want to be upfront about that rather than overstate the overlap. What transfers directly is the sensor-fusion architecture itself and the real-time embedded discipline around it: profiling compute and power on constrained hardware, validating fused estimates against ground truth, and shipping systems that have to work reliably in the field, not just in a lab benchmark.

I am a Canadian permanent resident, eligible to work in Canada without sponsorship, and I would relocate to Calgary for this role. I am available immediately and would welcome the chance to discuss how this background applies to Hexagon Autonomous Solutions' sensor fusion work.

Sincerely, \\
Nishal Stanislaus Silva

\end{document}
